% =============================================================================
% FORMALIA-MASCHINERIE -- hier musst du nichts aendern
% =============================================================================
% Diese Datei wird als erstes nach \documentclass geladen. Sie muss VOR
% fontspec, geometry und biblatex laufen, weil Schrift, Raender und Zitierstil
% Ladeoptionen dieser Pakete sind und sich nachtraeglich nicht mehr aendern
% lassen. Genau daran ist ein frueherer Entwurf gescheitert, der die Schalter
% in skripte/meta.tex legte -- die wird erst kurz vor \begin{document} gelesen.
%
% Reihenfolge: Standardwerte -> Leitfadenprofil -> Overlay -> eigene Vorgaben.
% Was spaeter kommt, gewinnt.
% =============================================================================

% =============================================================================
% FORMALIA-SCHALTER -- hier stellst du deinen Leitfaden ein
% =============================================================================
% An der FOM gibt es nicht DEN einen Leitfaden. Welcher fuer dich gilt, steht
% in deiner Modulbeschreibung bzw. sagen dir deine Pruefenden. Stelle hier ein,
% welcher gilt -- die Vorlage setzt Raender, Schrift, Zitierstil und die
% leitfadenspezifischen Sonderregeln dann automatisch richtig.
%
% Hintergrund zu jeder Einstellung: ../anleitung/02_formalia.md
%
% WICHTIG: Individuelle Vorgaben deiner Pruefenden haben IMMER Vorrang vor dem
% Leitfaden. Solche Abweichungen traegst du nicht hier ein, sondern in
% formalia/profil_eigen.tex -- dann bleibt nachvollziehbar, was Leitfaden ist
% und was Pruefervorgabe.
% =============================================================================

% --- Welcher Leitfaden gilt? -------------------------------------------------
%   wi   = Gestaltung wissenschaftlicher Arbeiten, V1.4 (03/2024)
%          Wirtschaftsinformatik & Ingenieurwesen
%   jks  = Formale Gestaltung, Jaeger/Kuempel/Seng (01/2024)
%          allgemeiner FOM-Standard (Wirtschaft & Recht u. a.)
\def\FormaliaLeitfaden{jks}

% --- Zitierstil --------------------------------------------------------------
%   fussnote    = Kurzbeleg in der Fussnote  (JKS-Standard, WI zulaessig)
%   authoryear  = Harvard/Autor-Jahr im Text (JKS zulaessig, WI zulaessig)
%   ieee        = numerisch [12]             (nur WI, in der Informatik ueblich)
%
% ACHTUNG: Diesen Wert einmal am Anfang festlegen und dann NICHT mehr wechseln.
% Der Schalter aendert nur die Ausgabe, nicht den Sinn der bereits im Text
% gesetzten \autocite-Befehle. Ein spaeter Wechsel von fussnote auf authoryear
% verschiebt saemtliche Belege aus den Fussnoten in den Fliesstext -- das musst
% du dann sprachlich nacharbeiten.
\def\FormaliaZitierstil{fussnote}

% --- Schrift -----------------------------------------------------------------
%   tnr    = Times New Roman 12 pt (Standard beider Leitfaeden)
%            Gesetzt wird TeX Gyre Termes, ein metrisch identischer freier Klon.
%   arial  = Arial als zulaessige Alternative (WI 11 pt / JKS 11,5 pt)
%            Gesetzt wird TeX Gyre Heros, ein metrisch identischer freier Klon.
\def\FormaliaSchrift{tnr}

% --- Optionales Overlay ------------------------------------------------------
%   ifes   = zusaetzliche Vorgaben fuer empirische Arbeiten (Gansser u. a.)
%   kein   = kein Overlay
% Das Overlay ersetzt den Leitfaden nicht, es ergaenzt ihn.
\def\FormaliaOverlay{kein}


% --- Standardwerte (gelten, bis ein Profil sie ueberschreibt) ----------------
\def\formaliaLeitfadenName{--}
\def\formaliaLeitfadenStand{--}
\def\formaliaGeprueftAm{--}

\def\formaliaRandLinks{4cm}
\def\formaliaRandRechts{2cm}
\def\formaliaRandOben{4cm}
\def\formaliaRandUnten{2cm}

\def\formaliaSchriftHaupt{TeX Gyre Termes}
\def\formaliaSchriftGroesse{12pt}

\def\formaliaSeitenzahlPos{C}          % C = oben mittig, R = oben rechts
\def\formaliaBibOpt{style=ext-authoryear-ibid,autocite=footnote,%
  mergedate=false,innamebeforetitle,dashed=false}

\newif\ifformaliaQuelleEigen           % "Quelle: Eigene Darstellung" verlangt?
\newif\ifformaliaSperrvermerkImTOC     % Sperrvermerk im Inhaltsverzeichnis?
\newif\ifformaliaInternetquellenSeparat

% --- Leitfadenprofil laden ---------------------------------------------------
% Bewusst als explizite Fallunterscheidung statt \input{...\FormaliaLeitfaden}:
% lesbarer, und ein Tippfehler im Schalter liefert eine klare Fehlermeldung
% statt "File not found".
\def\formaliaInternwi{wi}
\def\formaliaInternjks{jks}
\ifx\FormaliaLeitfaden\formaliaInternwi
  % =============================================================================
% PROFIL: "Gestaltung wissenschaftlicher Arbeiten", V1.4, Maerz 2024.
%         Wirtschaftsinformatik & Ingenieurwesen.
% -----------------------------------------------------------------------------
% Leitfaden-Fassung : V1.4, Maerz 2024
% Ausgewertet am    : 2026-08-26
% Belege            : ../anleitung/02_formalia.md
%
% Leitfaeden werden ueberarbeitet. Hol dir die aktuelle Fassung aus dem
% FOM Online-Campus und pruefe die Werte hier dagegen, bevor du abgibst.
% =============================================================================

\def\formaliaLeitfadenName{dem FOM-Leitfaden Wirtschaftsinformatik}
\def\formaliaLeitfadenStand{V1.4, Maerz 2024}
\def\formaliaGeprueftAm{2026-08-26}

% WI nennt die Masse "ungefaehr". Die Vorlage trifft die strengere Lesart.
\def\formaliaRandLinks{4cm}
\def\formaliaRandRechts{2cm}
\def\formaliaRandOben{4cm}
\def\formaliaRandUnten{2cm}

\def\formaliaSeitenzahlPos{C}

% WI verbietet den Vermerk "Quelle: Eigene Darstellung" ausdruecklich.
\formaliaQuelleEigenfalse

% WI: Sperrvermerk ohne Seitenzahl und NICHT im Inhaltsverzeichnis.
\formaliaSperrvermerkImTOCfalse

% WI trennt Internetquellen nicht ab.
\formaliaInternetquellenSeparatfalse

\else
  \ifx\FormaliaLeitfaden\formaliaInternjks
    % =============================================================================
% PROFIL: Jaeger/Kuempel/Seng -- "Formale Gestaltung wissenschaftlicher
%         Arbeiten", Januar 2024. Allgemeiner FOM-Standard.
% -----------------------------------------------------------------------------
% Leitfaden-Fassung : Januar 2024
% Ausgewertet am    : 2026-08-26
% Belege            : ../anleitung/02_formalia.md
%
% Leitfaeden werden ueberarbeitet. Hol dir die aktuelle Fassung aus dem
% FOM Online-Campus und pruefe die Werte hier dagegen, bevor du abgibst.
% =============================================================================

\def\formaliaLeitfadenName{dem FOM-Leitfaden von Jaeger, Kuempel und Seng}
\def\formaliaLeitfadenStand{Januar 2024}
\def\formaliaGeprueftAm{2026-08-26}

% Seitenlayout: der breite linke Rand ist Korrektur- und Bindungsrand.
\def\formaliaRandLinks{4cm}
\def\formaliaRandRechts{2cm}
\def\formaliaRandOben{4cm}
\def\formaliaRandUnten{2cm}

% Seitenzahl oben. JKS laesst mittig und "oben rechts am Textrand" zu.
\def\formaliaSeitenzahlPos{C}

% JKS verlangt bei selbst erstellten Abbildungen/Tabellen den Vermerk
% "Quelle: Eigene Darstellung". (WI verbietet ihn -- echter Widerspruch.)
\formaliaQuelleEigentrue

% JKS: Sperrvermerk mit roemischer Seitenzahl UND im Inhaltsverzeichnis.
\formaliaSperrvermerkImTOCtrue

% JKS: Internetquellen kommen separat ans Ende des Literaturverzeichnisses.
\formaliaInternetquellenSeparattrue

  \else
    \errmessage{[formalia] Unbekannter Leitfaden in formalia/konfig.tex.
      Erlaubt sind: wi, jks}
  \fi
\fi

% --- Zitierstil --------------------------------------------------------------
\def\formaliaInternfussnote{fussnote}
\def\formaliaInternauthoryear{authoryear}
\def\formaliaInternieee{ieee}
\ifx\FormaliaZitierstil\formaliaInternfussnote
  \def\formaliaBibOpt{style=ext-authoryear-ibid,autocite=footnote,%
    mergedate=false,innamebeforetitle,dashed=false}
\else
  \ifx\FormaliaZitierstil\formaliaInternauthoryear
    \def\formaliaBibOpt{style=authoryear,autocite=inline,%
      mergedate=false,dashed=false}
  \else
    \ifx\FormaliaZitierstil\formaliaInternieee
      \ifx\FormaliaLeitfaden\formaliaInternjks
        \errmessage{[formalia] Zitierstil ieee ist nach JKS nicht zulaessig.
          JKS erlaubt Chicago in Fussnoten oder Harvard im Text.}
      \fi
      \def\formaliaBibOpt{style=numeric-comp,autocite=inline,sorting=none}
    \else
      \errmessage{[formalia] Unbekannter Zitierstil in formalia/konfig.tex.
        Erlaubt sind: fussnote, authoryear, ieee}
    \fi
  \fi
\fi

% --- Schrift -----------------------------------------------------------------
\def\formaliaInterntnr{tnr}
\def\formaliaInternarial{arial}
\ifx\FormaliaSchrift\formaliaInterntnr
  \def\formaliaSchriftHaupt{TeX Gyre Termes}
  \def\formaliaSchriftGroesse{12pt}
\else
  \ifx\FormaliaSchrift\formaliaInternarial
    \def\formaliaSchriftHaupt{TeX Gyre Heros}
    % WI 11 pt, JKS 11,5 pt -- KOMA kennt keine halben Punkt als Klassenoption,
    % deshalb wird 11.5pt ueber \KOMAoptions gesetzt.
    \ifx\FormaliaLeitfaden\formaliaInternwi
      \def\formaliaSchriftGroesse{11pt}
    \else
      \def\formaliaSchriftGroesse{11.5pt}
    \fi
  \else
    \errmessage{[formalia] Unbekannte Schrift in formalia/konfig.tex.
      Erlaubt sind: tnr, arial}
  \fi
\fi

% --- Overlay -----------------------------------------------------------------
\def\formaliaInternifes{ifes}
\ifx\FormaliaOverlay\formaliaInternifes
  % =============================================================================
% OVERLAY: ifes -- Leitfaden fuer empirische Arbeiten (Gansser u. a.)
% -----------------------------------------------------------------------------
% Ein Overlay ersetzt den Leitfaden NICHT, es ergaenzt ihn. Geladen wird es
% zusaetzlich zu profil_wi bzw. profil_jks.
%
% Ausgewertet am : 2026-08-26
% Beleg          : ../anleitung/02_formalia.md
% =============================================================================

% ifes laesst den rechten Rand bis auf 1 cm zu. Die Vorlage bleibt bewusst bei
% 2 cm -- schmaler ist zulaessig, aber nicht besser lesbar. Wer den Platz
% braucht, setzt in profil_eigen.tex \def\formaliaRandRechts{1cm}.

% Inhaltliche ifes-Vorgabe, die LaTeX nicht pruefen kann und die deshalb in
% deiner FORMALIA.md stehen muss:
%   Der empirische Teil muss mehr als 50 % der Arbeit ausmachen.

\fi

% --- Eigene Vorgaben (Pruefende schlagen Leitfaden) --------------------------
\InputIfFileExists{formalia/profil_eigen}{}{}

% --- Werte an die Pakete uebergeben ------------------------------------------
% \edef expandiert die Makros JETZT, damit die Pakete fertige Werte sehen.
\edef\formaliaInterngeo{a4paper,%
  left=\formaliaRandLinks,right=\formaliaRandRechts,%
  top=\formaliaRandOben,bottom=\formaliaRandUnten,%
  headheight=14pt,headsep=1cm}
\expandafter\PassOptionsToPackage\expandafter{\formaliaInterngeo}{geometry}
\expandafter\PassOptionsToPackage\expandafter{\formaliaBibOpt}{biblatex}

\edef\formaliaInternfs{\formaliaSchriftGroesse}
\KOMAoptions{fontsize=\formaliaInternfs}

% --- Helfer fuer den Fliesstext ----------------------------------------------
% Unter selbst erstellten Abbildungen/Tabellen verwenden. JKS verlangt den
% Vermerk, WI verbietet ihn ausdruecklich -- dieser Befehl macht das Richtige.
\newcommand{\formaliaQuellenvermerk}{%
  \ifformaliaQuelleEigen Quelle: Eigene Darstellung\fi}

% Fuer die Methodik-Passage der Einleitung: den gewaehlten Leitfaden nennen
% ist selbst eine Formvorgabe.
\newcommand{\formaliaLeitfadenSatz}{%
  Die formale Gestaltung folgt \formaliaLeitfadenName\ (\formaliaLeitfadenStand).}
